\documentclass[10pt]{extarticle}
\usepackage[a4paper,margin=1.5cm]{geometry}
\usepackage[T1]{fontenc}
\usepackage{amsmath,amssymb}
\usepackage{graphicx}
\usepackage[export]{adjustbox}
\usepackage{multicol}
\usepackage{enumitem}
\usepackage{booktabs,array,tabularx}
\setlength{\columnsep}{0.7cm}
\setlength{\parindent}{0pt}
\setlength{\parskip}{3pt}
\pagestyle{empty}
\makeatletter
% figure registry (figdefs.tex is generated by the extract step)
\newcommand{\qfigdef}[4]{\@namedef{qf@file@#1}{#2}\@namedef{qf@w@#1}{#3}\@namedef{qf@wide@#1}{#4}}
\newcommand{\qf@check}[1]{\@ifundefined{qf@file@#1}{\PackageError{score2}{Unknown figure id #1}{}}{}}
\newcommand{\qf@img}[4][]{\includegraphics[width=#2\dimexpr\@nameuse{qf@w@#4}pt\relax,max width=#3,max height=0.45\textheight,#1]{\@nameuse{qf@file@#4}}}
% Full-width questions: finalize wraps every question that contains a wide figure in \qfullbegin...\qfullend,
% so the WHOLE question (stem, figures, options) is set across the page in reading order, then the two columns
% resume. The question's head - its number, stem and first figure (\qfullsplit marks the end of the head) -
% is kept on one page, so a number is never left at a page foot with its figure overleaf; the rest may break
% across pages normally, so no big gaps are left. Without a \qfullsplit, the whole question is kept together
% if it is under 45% of a page. Never write these yourself.
\newif\ifqf@full
\newif\ifqf@boxed
\newsavebox{\qf@box}
\newcommand{\qf@open}{\global\qf@fulltrue\par\global\qf@boxedtrue
  \setbox\qf@box\vbox\bgroup\hsize\textwidth\linewidth\textwidth}
\newcommand{\qf@emit}[1]{\par\egroup\global\qf@boxedfalse
  \ifdim\dimexpr\ht\qf@box+\dp\qf@box\relax<#1\textheight\noindent\usebox{\qf@box}\par\else\unvbox\qf@box\fi}
\newcommand{\qfullsplit}{\ifqf@boxed\qf@emit{0.85}\fi}
\newcommand{\qf@close}{\ifqf@boxed\qf@emit{0.45}\fi\global\qf@fullfalse}
\newcommand{\qfullbegin}{\end{multicols}\qf@open}
\newcommand{\qfullend}{\qf@close\begin{multicols}{2}}
\newcommand{\qfullstart}{\qf@open}
\newcommand{\qfulldone}{\qf@close}
% \qfig[scale]{id}: block figure, centred, directly where it is written (size set by the tool).
% Fallback if a wide figure is not inside a full-width question: break out of the columns in place - no floats.
\newcommand{\qfig}[2][1]{\qf@check{#2}%
  \ifqf@full\qf@blk{#1}{#2}\else
  \ifnum\@nameuse{qf@wide@#2}=1 \expandafter\@firstoftwo\else\expandafter\@secondoftwo\fi
  {\end{multicols}\begin{center}\qf@img{#1}{\textwidth}{#2}\end{center}\begin{multicols}{2}}%
  {\qf@blk{#1}{#2}}\fi}
\newcommand{\qf@blk}[2]{\par\smallskip{\centering\qf@img{#1}{\linewidth}{#2}\par}\smallskip}
% \qfigi[scale]{id}: inline figure (inside option lists / tables / minipages); never breaks columns.
\newcommand{\qfigi}[2][1]{\qf@check{#2}\qf@img[valign=c]{#1}{\linewidth}{#2}}
\makeatother
% structure macros - lint reads these to check completeness and figure placement
\newcommand{\papersec}[1]{\par\bigskip{\centering\Large\bfseries PAPER-#1\par}\medskip}
\newcommand{\partsec}[2]{\par\medskip\noindent{\large\bfseries PART-#1: #2}\par\smallskip}
\newcommand{\sect}[1]{\par\smallskip\noindent\textbf{#1}\par}
\newcommand{\qn}[1]{\par\smallskip\noindent\textbf{#1.}~}
\newcommand{\paraq}[2]{\par\smallskip\noindent\textbf{Paragraph for Questions #1 and #2}\par}
\newenvironment{qoptions}{\begin{itemize}[leftmargin=1.8em,itemsep=1pt,topsep=2pt,label={}]}{\end{itemize}}
\newcommand{\opt}[2]{\item[(#1)] #2}
% figure registry (generated by score2_tools.py finalize)
\qfigdef{p1_phy_s1_q01_1}{Enthusiast Pre-Score - Test 01 (04 Nov 2020)-img/p1_phy_s1_q01_1.png}{206.2}{0}
\qfigdef{p1_phy_s1_q02_1}{Enthusiast Pre-Score - Test 01 (04 Nov 2020)-img/p1_phy_s1_q02_1.png}{157.5}{0}
\qfigdef{p1_phy_s1_q03_1}{Enthusiast Pre-Score - Test 01 (04 Nov 2020)-img/p1_phy_s1_q03_1.png}{139.6}{0}
\qfigdef{p1_phy_s1_q03_A}{Enthusiast Pre-Score - Test 01 (04 Nov 2020)-img/p1_phy_s1_q03_A.png}{126.4}{0}
\qfigdef{p1_phy_s1_q03_B}{Enthusiast Pre-Score - Test 01 (04 Nov 2020)-img/p1_phy_s1_q03_B.png}{126.4}{0}
\qfigdef{p1_phy_s1_q03_C}{Enthusiast Pre-Score - Test 01 (04 Nov 2020)-img/p1_phy_s1_q03_C.png}{127.0}{0}
\qfigdef{p1_phy_s1_q03_D}{Enthusiast Pre-Score - Test 01 (04 Nov 2020)-img/p1_phy_s1_q03_D.png}{127.0}{0}
\qfigdef{p1_phy_s1_q06_1}{Enthusiast Pre-Score - Test 01 (04 Nov 2020)-img/p1_phy_s1_q06_1.png}{118.0}{0}
\qfigdef{p1_phy_s2_q07_1}{Enthusiast Pre-Score - Test 01 (04 Nov 2020)-img/p1_phy_s2_q07_1.png}{237.8}{0}
\qfigdef{p1_phy_s2_q08_1}{Enthusiast Pre-Score - Test 01 (04 Nov 2020)-img/p1_phy_s2_q08_1.png}{114.4}{0}
\qfigdef{p1_phy_s2_q09_1}{Enthusiast Pre-Score - Test 01 (04 Nov 2020)-img/p1_phy_s2_q09_1.png}{126.5}{0}
\qfigdef{p1_phy_s2_q10_1}{Enthusiast Pre-Score - Test 01 (04 Nov 2020)-img/p1_phy_s2_q10_1.png}{137.3}{0}
\qfigdef{p1_phy_s3_q01_1}{Enthusiast Pre-Score - Test 01 (04 Nov 2020)-img/p1_phy_s3_q01_1.png}{139.6}{0}
\qfigdef{p1_phy_s3_q02_1}{Enthusiast Pre-Score - Test 01 (04 Nov 2020)-img/p1_phy_s3_q02_1.png}{124.8}{0}
\qfigdef{p1_phy_s3_q03_1}{Enthusiast Pre-Score - Test 01 (04 Nov 2020)-img/p1_phy_s3_q03_1.png}{98.4}{0}
\qfigdef{p1_phy_s3_q04_1}{Enthusiast Pre-Score - Test 01 (04 Nov 2020)-img/p1_phy_s3_q04_1.png}{75.5}{0}
\qfigdef{p1_phy_s3_q05_1}{Enthusiast Pre-Score - Test 01 (04 Nov 2020)-img/p1_phy_s3_q05_1.png}{173.0}{0}
\qfigdef{p1_phy_s3_q06_1}{Enthusiast Pre-Score - Test 01 (04 Nov 2020)-img/p1_phy_s3_q06_1.png}{83.8}{0}
\qfigdef{p1_chem_s1_q04_A}{Enthusiast Pre-Score - Test 01 (04 Nov 2020)-img/p1_chem_s1_q04_A.png}{105.4}{0}
\qfigdef{p1_chem_s1_q04_B}{Enthusiast Pre-Score - Test 01 (04 Nov 2020)-img/p1_chem_s1_q04_B.png}{105.4}{0}
\qfigdef{p1_chem_s1_q04_C}{Enthusiast Pre-Score - Test 01 (04 Nov 2020)-img/p1_chem_s1_q04_C.png}{105.4}{0}
\qfigdef{p1_chem_s1_q04_D}{Enthusiast Pre-Score - Test 01 (04 Nov 2020)-img/p1_chem_s1_q04_D.png}{105.4}{0}
\qfigdef{p1_chem_s2_q07_A}{Enthusiast Pre-Score - Test 01 (04 Nov 2020)-img/p1_chem_s2_q07_A.png}{29.3}{0}
\qfigdef{p1_chem_s2_q07_B}{Enthusiast Pre-Score - Test 01 (04 Nov 2020)-img/p1_chem_s2_q07_B.png}{61.1}{0}
\qfigdef{p1_chem_s2_q07_C}{Enthusiast Pre-Score - Test 01 (04 Nov 2020)-img/p1_chem_s2_q07_C.png}{40.2}{0}
\qfigdef{p1_chem_s2_q07_D}{Enthusiast Pre-Score - Test 01 (04 Nov 2020)-img/p1_chem_s2_q07_D.png}{72.5}{0}
\qfigdef{p1_chem_s2_q08_1}{Enthusiast Pre-Score - Test 01 (04 Nov 2020)-img/p1_chem_s2_q08_1.png}{215.5}{0}
\qfigdef{p1_chem_s2_q08_2}{Enthusiast Pre-Score - Test 01 (04 Nov 2020)-img/p1_chem_s2_q08_2.png}{29.0}{0}
\qfigdef{p1_chem_s2_q08_3}{Enthusiast Pre-Score - Test 01 (04 Nov 2020)-img/p1_chem_s2_q08_3.png}{31.6}{0}
\qfigdef{p1_chem_s3_q02_1}{Enthusiast Pre-Score - Test 01 (04 Nov 2020)-img/p1_chem_s3_q02_1.png}{154.0}{0}
\qfigdef{p1_chem_s3_q05_1}{Enthusiast Pre-Score - Test 01 (04 Nov 2020)-img/p1_chem_s3_q05_1.png}{380.7}{1}
\qfigdef{p1_chem_s3_q06_1}{Enthusiast Pre-Score - Test 01 (04 Nov 2020)-img/p1_chem_s3_q06_1.png}{112.8}{0}
\qfigdef{p2_phy_s1_q01_1}{Enthusiast Pre-Score - Test 01 (04 Nov 2020)-img/p2_phy_s1_q01_1.png}{264.6}{1}
\qfigdef{p2_phy_s1_q02_1}{Enthusiast Pre-Score - Test 01 (04 Nov 2020)-img/p2_phy_s1_q02_1.png}{190.3}{0}
\qfigdef{p2_phy_s1_q03_2}{Enthusiast Pre-Score - Test 01 (04 Nov 2020)-img/p2_phy_s1_q03_2.png}{126.4}{0}
\qfigdef{p2_phy_s1_q03_1}{Enthusiast Pre-Score - Test 01 (04 Nov 2020)-img/p2_phy_s1_q03_1.png}{137.0}{0}
\qfigdef{p2_phy_s1_q04_1}{Enthusiast Pre-Score - Test 01 (04 Nov 2020)-img/p2_phy_s1_q04_1.png}{277.1}{1}
\qfigdef{p2_phy_s1_q05_1}{Enthusiast Pre-Score - Test 01 (04 Nov 2020)-img/p2_phy_s1_q05_1.png}{188.5}{0}
\qfigdef{p2_phy_s1_q06_1}{Enthusiast Pre-Score - Test 01 (04 Nov 2020)-img/p2_phy_s1_q06_1.png}{158.8}{0}
\qfigdef{p2_phy_s2_q02_1}{Enthusiast Pre-Score - Test 01 (04 Nov 2020)-img/p2_phy_s2_q02_1.png}{110.6}{0}
\qfigdef{p2_phy_s2_q03_1}{Enthusiast Pre-Score - Test 01 (04 Nov 2020)-img/p2_phy_s2_q03_1.png}{62.0}{0}
\qfigdef{p2_phy_s2_q04_1}{Enthusiast Pre-Score - Test 01 (04 Nov 2020)-img/p2_phy_s2_q04_1.png}{75.9}{0}
\qfigdef{p2_phy_s2_q05_1}{Enthusiast Pre-Score - Test 01 (04 Nov 2020)-img/p2_phy_s2_q05_1.png}{138.1}{0}
\qfigdef{p2_phy_s2_q06_1}{Enthusiast Pre-Score - Test 01 (04 Nov 2020)-img/p2_phy_s2_q06_1.png}{156.5}{0}
\qfigdef{p2_phy_s3_q01_1}{Enthusiast Pre-Score - Test 01 (04 Nov 2020)-img/p2_phy_s3_q01_1.png}{94.3}{0}
\qfigdef{p2_phy_s3_q02_1}{Enthusiast Pre-Score - Test 01 (04 Nov 2020)-img/p2_phy_s3_q02_1.png}{103.5}{0}
\qfigdef{p2_phy_s3_q03_1}{Enthusiast Pre-Score - Test 01 (04 Nov 2020)-img/p2_phy_s3_q03_1.png}{175.5}{0}
\qfigdef{p2_phy_s3_q04_1}{Enthusiast Pre-Score - Test 01 (04 Nov 2020)-img/p2_phy_s3_q04_1.png}{104.1}{0}
\qfigdef{p2_phy_s3_q05_1}{Enthusiast Pre-Score - Test 01 (04 Nov 2020)-img/p2_phy_s3_q05_1.png}{164.5}{0}
\qfigdef{p2_phy_s3_q06_1}{Enthusiast Pre-Score - Test 01 (04 Nov 2020)-img/p2_phy_s3_q06_1.png}{237.4}{0}
\qfigdef{p2_chem_s1_q05_1}{Enthusiast Pre-Score - Test 01 (04 Nov 2020)-img/p2_chem_s1_q05_1.png}{148.2}{0}
\qfigdef{p2_chem_s2_q05_1}{Enthusiast Pre-Score - Test 01 (04 Nov 2020)-img/p2_chem_s2_q05_1.png}{363.1}{1}
\qfigdef{p2_chem_s3_q03_1}{Enthusiast Pre-Score - Test 01 (04 Nov 2020)-img/p2_chem_s3_q03_1.png}{387.1}{1}
\qfigdef{p2_chem_s3_q04_1}{Enthusiast Pre-Score - Test 01 (04 Nov 2020)-img/p2_chem_s3_q04_1.png}{360.2}{1}
\raggedcolumns
\makeatletter\@beginparpenalty=10000\makeatother
\newcolumntype{L}[1]{>{\raggedright\arraybackslash}p{#1}}
\begin{document}
\begin{center}{\Large\bfseries Enthusiast Pre-Score --- Test 01 (04 Nov 2020)}\end{center}
\begin{multicols}{2}
\raggedcolumns
\papersec{1}
\partsec{1}{PHYSICS}
\sect{SECTION-I(i)}
\qn{1} Two blocks and a pulley are connected by inextensible cords as shown. The relative velocity of block A with respect to block B is 0.8 m/s to the left at time t = 0 and 0.4 m/s to the left when t = 0.25 s. Knowing that the angular acceleration of the pulley is constant, the relative acceleration of block A with respect to block B is :-\par
\qfig{p1_phy_s1_q01_1}
\begin{qoptions}
\opt{A}{1.6 m/s$^2$, left}
\opt{B}{1.6 m/s$^2$, right}
\opt{C}{0.8 m/s$^2$, left}
\opt{D}{3.2 m/s$^2$, right}
\end{qoptions}
\qn{2} A uniform ring of radius R and mass m is fitted with a massless rod AB along its diameter as shown in figure. The coefficient of friction between ring and surface is $\mu$. When the system is released from rest, the ring moves such that rod AB remains vertical. The value of r is :-\par
\qfig{p1_phy_s1_q02_1}
\begin{qoptions}
\opt{A}{$R\left(1-\dfrac{3\mu}{2(1+\mu)}\right)$}
\opt{B}{$R\left(1-\dfrac{\mu}{2(1+\mu)}\right)$}
\opt{C}{$R\left(2-\dfrac{3\mu}{2(1+\mu)}\right)$}
\opt{D}{$R\left(1-\dfrac{3\mu}{(1+\mu)}\right)$}
\end{qoptions}
\qn{3} A disc of mass m and radius R is under pure rolling motion with constant velocity v$_0$ on a smooth surface. The centre of mass of the disc C is offset from centre O by a distance r. A time varying couple M is applied such that disc continues to roll with a constant velocity. The correct plot of variation of couple M with time t will be [Clockwise torque is +ve]\par
\qfig{p1_phy_s1_q03_1}
\begin{qoptions}
\opt{A}{\qfigi{p1_phy_s1_q03_A}}
\opt{B}{\qfigi{p1_phy_s1_q03_B}}
\opt{C}{\qfigi{p1_phy_s1_q03_C}}
\opt{D}{\qfigi{p1_phy_s1_q03_D}}
\end{qoptions}
\qn{4} A particle moves along a circle of radius R, starting from rest, where its velocity varies directly as the square root of the distance `s' described. The angle between the resultant acceleration and the velocity vector, will be :
\begin{qoptions}
\opt{A}{$\tan^{-1}\left(\dfrac sR\right)$}
\opt{B}{$\tan^{-1}\left(\dfrac{2s}R\right)$}
\opt{C}{$\tan^{-1}\left(\dfrac s{2R}\right)$}
\opt{D}{$\tan^{-1}\left(\dfrac s{4R}\right)$}
\end{qoptions}
\qn{5} A smooth sphere is moving on a horizontal surface with velocity vector $2\hat i+2\hat j$ immediately before it hits a vertical wall. The wall is parallel to $\hat j$ vector and the coefficient of restitution between the sphere and the wall is $e=\dfrac12$. The velocity vector of the sphere after it hits the wall is
\begin{qoptions}
\opt{A}{$\hat i-\hat j$}
\opt{B}{$-\hat i+2\hat j$}
\opt{C}{$-\hat i-\hat j$}
\opt{D}{$2\hat i-\hat j$}
\end{qoptions}
\qn{6} A block of mass 1 kg is resting on a rough horizontal surface with coefficient of friction 0.2. If an impulse $\left(10\hat i-10\hat j\right)$ N-s is given to block, then the distance travelled by the block till it stops (after the impulse) is\par
\qfig{p1_phy_s1_q06_1}
\begin{qoptions}
\opt{A}{10 m}
\opt{B}{0 m}
\opt{C}{16 m}
\opt{D}{4m}
\end{qoptions}
\sect{SECTION-I(ii)}
\qn{7} A particle is projected from origin with speed u=50 m/s at an angle 30$^\circ$ with X-axis. Due to wind particle acquires an extra acceleration a = 6 m/s$^2$ in $-$ve x-direction. Path of particle is shown in figure.\par
\qfig{p1_phy_s2_q07_1}
x - Coordinates of point A and B are same\par
y- Co-ordinates of P and Q are same.\par
Point N on the X-axis, point M on the Y-axis.\par
Time taken by particle from O to P, O to Q, O to A and O to B are t$_1$, t$_2$, t$_3$ and t$_4$ respectively. (Assuming Y along vertical and their is no surface in the space) (g = 10 m/s$^2$ and constant).
\begin{qoptions}
\opt{A}{The value of t$_1$ + t$_2$ is 5 sec.}
\opt{B}{The value of t$_3$ + t$_4$ is 25 sec.}
\opt{C}{The length of ON is $125\sqrt3\,(5-\sqrt3)$ m}
\opt{D}{The length of OM is $\dfrac{625}{\sqrt3}\left(\dfrac5{\sqrt3}-1\right)$ m}
\end{qoptions}
\qn{8} A man of height h$_0$ = 2m and mass m = 50 kg is bungee jumping from a platform situated at a height h = 25 m above a lake. One end of an elastic rope is attached to his foot and the other end is fixed to the platform. He starts falling from rest in vertical position. The length and elastic properties of the rope are chosen so that his speed will have been reduced to zero at instant when his head reaches the surface of water. Ultimately the jumper is hanging from the rope with his head 8m above the water. (take g = 10 m/s$^2$) (air resistance is very-very small, center of mass of man at mid-point)\par
\qfig{p1_phy_s2_q08_1}
\begin{qoptions}
\opt{A}{The maximum acceleration achieved during the motion is approximately 40 m/s$^2$.}
\opt{B}{Natural length of elastic rope is 13m.}
\opt{C}{When finally man comes to at rest, elongation in the rope is 2m.}
\opt{D}{Total work done by elastic rope is 250 J.}
\end{qoptions}
\qn{9} A particle of mass m is projected up from the bottom of an smooth inclined plane with initial velocity v$_0$ at angle 45$^\circ$ with an inclined plane of inclination 30$^\circ$ as shown in figure. At the same time, a small block of same mass m is released from rest at a height h. The particle hits the block at some point on inclined plane. (take g = 10 m/s$^2$) (size of the block is negligible)\par
\qfig{p1_phy_s2_q09_1}
\begin{qoptions}
\opt{A}{The value of height h is $\dfrac{v_0^2}{\sqrt3\mathrm g}$\,.}
\opt{B}{The value of height h is $\dfrac{\sqrt3v_0^2}{\mathrm g}$\,.}
\opt{C}{If the particle stick to the block after collision, the velocity of block parallel to the inclined plane just after collision will be $\dfrac{v_0}{\sqrt2}$\,.}
\opt{D}{If the particle stick to the block after collision, the velocity of block parallel to the inclined plane just after collision will be $\dfrac{v_0}{2\sqrt2}\left(\dfrac4{\sqrt3}-1\right)$\,.}
\end{qoptions}
\qn{10} The hollow cylinder $A$ and pulley $B$ turn without friction about stationary horizontal axles that pass through their centres. A light rope wrapped around the cylinder, passes over the pulley and has a 3 kg box $C$ suspended from its free end. There is no slipping between the rope and the pulley surface. The hollow cylinder A has mass 5 kg and radius 40 cm. The pulley is a uniform disk with mass 2 kg and radius 20 cm. The box is released from rest and descends as the rope unwraps from the cylinder.\par
\qfig{p1_phy_s2_q10_1}
\begin{qoptions}
\opt{A}{The speed of the box when it has fallen 1.50 m is $\sqrt{10}$ m/s.}
\opt{B}{The speed of the box when it has fallen 1.50 m is 6.28 m/s.}
\opt{C}{The ratio of angular speed of cylinder and disk is 1 : 2.}
\opt{D}{The ratio of angular speed of cylinder and disk is 2 : 3.}
\end{qoptions}
\qn{11} A thin uniform rod of mass m and length $\ell$ is constrained to rotate about its upper end. When it is at rest, it receives an impulse $J$ at its lower most point, normal to its length. Immediately after impact :-
\begin{qoptions}
\opt{A}{The angular momentum of the rod about its upper end is $J\ell$}
\opt{B}{The angular momentum of the rod about its upper end is $\dfrac{3J}{m\ell}$}
\opt{C}{The kinetic energy of the rod is $\dfrac{3J^2}{2m}$}
\opt{D}{The linear velocity of the midpoint of the rod is $\dfrac{3J}{2m}$}
\end{qoptions}
\qn{12} A particle of mass m is acted upon by a conservative force only and moves along path given by $x=x_0\cos(\omega t)$ and $y=y_0\sin(\omega t)$ where $x_0\ y_0$ and $\omega$ are constant. Take potential energy U = 0 at x = 0 \& y = 0. Mark the \textbf{CORRECT} option :-
\begin{qoptions}
\opt{A}{Force acting on the particle is $\mathrm F=\mathrm m\omega^2\,\mathrm x(-\hat i)+\mathrm m\omega^2\mathrm y(-\hat j)$}
\opt{B}{Potential energy at any x \& y is given by $\mathrm U=\dfrac12\mathrm m\ \omega^2(\mathrm x^2+\mathrm y^2)$}
\opt{C}{Total energy of the particle will remain constant.}
\opt{D}{Kinetic energy of the particle at x = x$_0$, y = y$_0$ is $\dfrac12\mathrm m(\omega x_0)^2$}
\end{qoptions}
\sect{SECTION-II}
\qn{1} A solid cylinder of radius r = 5m is placed with no initial velocity on a moving belt as shown such that it touches a fixed vertical rod. Assuming that the coefficient of friction at all contact surfaces is $\mu$ = 0.5, determine the angular acceleration of the cylinder in rad/s$^2$. Take g = 10 m/s$^2$.\par
\qfig{p1_phy_s3_q01_1}
\qn{2} In the figure shown, a long cart is to be moved with constant velocity v = 5 m/s on a smooth horizontal surface. As it moves, sand is falling from a stationary hopper on to the cart with negligible velocity at constant rate 0.1 kg/s and is sticking to the cart. Obtain the power (in watts) supplied by external agent applying that force.\par
\qfig{p1_phy_s3_q02_1}
\qn{3} The sphere at P is given a downward velocity v$_0$ and swings in a vertical plane at the end of a rope of $\ell$ = 1m attached to a support at O. The rope breaks at angle 30$^\circ$ from horizontal, knowing that it can withstand a maximum tension equal to three times the weight of the sphere. Then the value of v$_0$ will be :\par
\qfig{p1_phy_s3_q03_1}
\qn{4} A rod of mass 10 kg and length 1m is at rest on plane horizontal smooth surface. A particle of same mass M strike one end with velocity 2 m/s perpendicular to rod, elastically. Find kinetic energy of upper half part of rod (in J) just after collision.\par
\qfig{p1_phy_s3_q04_1}
\qn{5} A rod of length `$\ell$' is standing vertically on smooth horizontal surface. It is disturbed slightly so that lower end turns towards step. Find angular speed (in rad/s) of rod just after hitting step (end of rod sticks to step after collision). Given : $\ell$ = 13 cm\par
\qfig{p1_phy_s3_q05_1}
\qn{6} Two blocks of masses m$_1$ = 5 kg and m$_2$ = 2 kg are connected by threads which pass over the pulleys as shown in the figure. The threads are massless and the pulleys are massless and smooth. The blocks can move only along the vertical direction. If the acceleration of m$_1$, if expressed in simplest form, is equal to $\left(\dfrac{2n}{n-1}\right)$ m/sec$^2$, calculate `n'. (Take g = 10 m/s$^2$)\par
\qfig{p1_phy_s3_q06_1}
\partsec{2}{CHEMISTRY}
\sect{SECTION-I(i)}
\qn{1} Which of the following is \textbf{CORRECT} statement?
\begin{qoptions}
\opt{A}{Concentrated HNO$_3$ give NO$_2$ gas when reacts with zinc.}
\opt{B}{Concentrated HNO$_3$ give HIO$_2$ when reacts with iodine}
\opt{C}{Dilute HNO$_3$ give NO$_2$ when reacts with copper.}
\opt{D}{Dilute HNO$_3$ give NH$_3$ when reacts with copper.}
\end{qoptions}
\qn{2} On burning hydrogen in air the colour of flame is?
\begin{qoptions}
\opt{A}{Green}
\opt{B}{Light bluish}
\opt{C}{Violet}
\opt{D}{Red}
\end{qoptions}
\qn{3} Which of the following gives Isocyanide test with CHCl$_3$ + KOH :
\begin{qoptions}
\opt{A}{N-methylcyclohexanamine}
\opt{B}{Acetamide}
\opt{C}{Aniline}
\opt{D}{N,N-dimethylaniline}
\end{qoptions}
\qn{4} The correct orientation of dipoles in pyrrole and pyridine is :\par
\begin{qoptions}
\opt{A}{\qfigi{p1_chem_s1_q04_A}}
\opt{B}{\qfigi{p1_chem_s1_q04_B}}
\opt{C}{\qfigi{p1_chem_s1_q04_C}}
\opt{D}{\qfigi{p1_chem_s1_q04_D}}
\end{qoptions}
\qn{5} For a dimerization reaction,\par
2 A(g) $\to$ A$_2$(g)\par
at 300 K, $\Delta U^0=-20$kJ mol$^{-1}$, $\Delta S^0=-30$ JK$^{-1}$ mol$^{-1}$\par
then the $\mathbf{lnK_{eq.}}$ for the dissociation of A$_2$ at 300 K will be .
\begin{qoptions}
\opt{A}{$-5.4$}
\opt{B}{$+5.4$}
\opt{C}{$-4.4$}
\opt{D}{$+4.4$}
\end{qoptions}
\qn{6} For the H$_2$O($l$) $\rightleftharpoons$ H$_2$O(g)\par
Select the incorrect option -
\begin{qoptions}
\opt{A}{$[G_{\mathrm{H_2O}(l)}]_{\mathrm{1atm,\,373K}}=[G_{\mathrm{H_2O(g)}}]_{\mathrm{1atm,\,373K}}$}
\opt{B}{$[G_{\mathrm{H_2O}(l)}]_{\mathrm{1atm,\,400K}}>[G_{\mathrm{H_2O(g)}}]_{\mathrm{1atm,\,400K}}$}
\opt{C}{$[G_{\mathrm{H_2O}(l)}]_{\mathrm{0.5atm,\,373K}}>[G_{\mathrm{H_2O(g)}}]_{\mathrm{0.5atm,\,373K}}$}
\opt{D}{$[G_{\mathrm{H_2O}(l)}]_{\mathrm{0.5atm,\,400K}}<[G_{\mathrm{H_2O(g)}}]_{\mathrm{0.5atm,\,400K}}$}
\end{qoptions}
\sect{SECTION-I(ii)}
\qn{7} Which of the following will react at rate faster than benzene during aromatic electrophilic substitution?
\begin{qoptions}
\opt{A}{\qfigi{p1_chem_s2_q07_A}}
\opt{B}{\qfigi{p1_chem_s2_q07_B}}
\opt{C}{\qfigi{p1_chem_s2_q07_C}}
\opt{D}{\qfigi{p1_chem_s2_q07_D}}
\end{qoptions}
\qn{8} Which are \textbf{CORRECT} statement(s) in the following?
\begin{qoptions}
\opt{A}{Rate of sulphonation of C$_6$H$_6$ and C$_6$D$_6$ are equal.}
\opt{B}{Rate of halogentation of\par\qfig[0.8]{p1_chem_s2_q08_1}}
\opt{C}{\qfigi{p1_chem_s2_q08_2} on reaction with Br$_2$ / H$_2$O forms tribromoderivative.}
\opt{D}{\qfigi{p1_chem_s2_q08_3} on reaction with Cl$_2$ at 500$^\circ$C will give electrophilic aromatic substitution reaction.}
\end{qoptions}
\qn{9} Which of the following reaction / properties is / are given by sulphuric acid?
\begin{qoptions}
\opt{A}{H$_2$SO$_4$ forms 2 series of salts and one of them is acidic salt}
\opt{B}{Solvolysis of salts of acids that are weaker than H$_2$SO$_4$\par KNO$_3$ + H$_2$SO$_4 \longrightarrow$ K$^{\oplus}$ + HSO$_4^{\ominus}$ + HNO$_3$}
\opt{C}{Protonation of compounds with lone pairs of electrons\par H$_2$O + H$_2$SO$_4 \longrightarrow$ H$_3$O$^{\oplus}$ + HSO$_4^{\ominus}$}
\opt{D}{Dehydration reactions\par H$_2$C$_2$O$_4$ + [conc. H$_2$SO$_4$] $\longrightarrow$ H$_2$O + CO + CO$_2$}
\end{qoptions}
\qn{10} Choose \textbf{CORRECT} statement(s) -
\begin{qoptions}
\opt{A}{Oxygen has three stable isotopes}
\opt{B}{Solubility of oxygen in salt water is less than normal fresh water}
\opt{C}{The blue colour of oxygen in the liquid and solid phase is due to electronic transition where O$_2$ molecules in the triplet ground state are excited to the singlet state.}
\opt{D}{Xenon trioxide is formed in the hydrolysis of XeF$_4$ and XeF$_6$}
\end{qoptions}
\qn{11} For N$_2$(g) + 3H$_2$(g) $\rightleftharpoons$ 2NH$_3$(g)\par
on increasing temperature at 1 atm which of the following is/are correct [Assume enthalpy of reaction is constant]
\begin{qoptions}
\opt{A}{Kp decreases}
\opt{B}{Kp increases}
\opt{C}{Unfavourable $\Delta S_{\mathrm{surr.}}$ increases}
\opt{D}{Favourable $\Delta S_{\mathrm{surr.}}$ decreases}
\end{qoptions}
\qn{12} For the process : H$_2$O(s) $\xrightarrow[(-10^\circ\mathrm C,1\mathrm{atm})]{}$ H$_2$O(g) $_{(473^\circ\mathrm C,1\mathrm{atm})}$\par
Given : $[(\mathrm{Cp})_{\mathrm{H_2O(s)}}=\dfrac{10\mathrm{cal}}{\mathrm{molK}}$, $\Delta H_{\mathrm{fusion\ (H_2O,\,0^\circ C)}}=2$ kcal/mol, $\Delta H_{\mathrm{vap(H_2O,100^\circ C)}}=10\dfrac{\mathrm{kcal}}{\mathrm{mol-K}}$,\par
$(\mathrm{Cp})_{\mathrm{H_2O(l)}}=20\dfrac{\mathrm{cal}}{\mathrm{mol-K}}$, $\ln\dfrac{273}{263}=0.04$, $\ln\dfrac{373}{273}=0.31$, $\ln2=0.7]$
\begin{qoptions}
\opt{A}{$\Delta H=17.1$ kcal/mol}
\opt{B}{$\Delta H=14.1$ kcal/mol}
\opt{C}{$\Delta S=46.33$ cal/K-mol}
\opt{D}{$\Delta S=20.2$ cal/K-mol}
\end{qoptions}
\sect{SECTION-II}
\qn{1} Tin is obtained from cassiterite by reduction with coke. Use the data given below to determine the minimum temperature (in K) at which the reduction of cassiterite by coke would take place as :\par
SnO$_2$(s) + C(s) $\to$ Sn(s) + CO$_2$(g)\par
At 298 K : $\Delta_f H^0(\mathrm{SnO_2(s)})=-584.0$ kJ mol$^{-1}$, $\Delta_f H^0(\mathrm{CO_2(g)})=-394.0$ kJ mol$^{-1}$\par
S$^0$(SnO$_2$(s)) = 56.0 J K$^{-1}$ mol$^{-1}$, S$^0$(Sn(s)) = 52.0 J K$^{-1}$ mol$^{-1}$ ,\par
S$^0$(C(s)) = 6.0 J K$^{-1}$ mol$^{-1}$ , S$^0$(CO$_2$(g)) = 210.0 J K$^{-1}$ mol$^{-1}$.\par
Assume that the enthalpies and the entropies are temperature independent.
\qn{2} Consider the reaction \textbf{A} $\rightleftharpoons$ \textbf{2B} at equilibrium at 1000 K. At time t', the temperature of the system was increased to 3000 K and the system was allowed to reach equilibrium. Throughout this experiment the partial pressure of \textbf{A} was maintained at 1 bar. Given below is the plot of the partial pressure of B with time. What is the ratio of the standard Gibbs energy of the reaction at 1000 K to that at 3000 K?\par
\qfig{p1_chem_s3_q02_1}
\qn{3} Consider the following molecules\par
IF$_3$, IF$_5$, IF, CO, CO$_2$, HCN, IF$_7$, HNC\par
if x is total number of molecule in which only unpaired electron(s) in the ground state of the central atom is participating in the bond formation.\par
y is total number of molecule in which coordinate bond is present\par
then calculate value of $\left(\dfrac xy\right)$
\qn{4} Consider the O$_2^{\oplus}$, O$_2$, O$_2^-$, O$_2^{2-}$ calculate total number of $\pi$-antibonding electron in species which have highest IR vibrational frequency?
\qfullbegin
\qn{5} Identify L.S. of compounds give prussian blue with FeSO$_4$?\par
\qfig{p1_chem_s3_q05_1}
\qfullend
\qn{6} Identify compounds give Biurate test?\par
(i) Cellulose \quad (ii) Protien \quad (iii) Nylon-6 \quad (iv) Starch\par
(v)\par\qfig[0.8]{p1_chem_s3_q06_1}\par(vi) H$_2$N$-$gly$-$Ala$-$Cys$-$COOH
\partsec{3}{MATHEMATICS}
\sect{SECTION-I(i)}
\qn{1} For equation $|x^2-3x-4|=|x-2|-1$ which of following is correct-
\begin{qoptions}
\opt{A}{Total number of solutions is 4}
\opt{B}{Total number of solutions is 3}
\opt{C}{Total number of rational solutions is 3}
\opt{D}{Sum of all solution is $\sqrt{16}+\sqrt8+\sqrt5-\sqrt6$}
\end{qoptions}
\qn{2} If $x=\log2$ and $y=\log3$, then $a+bx+cy=[\log1+\log(1+3)+\log(1+3+5)+.....+\log(1+3+5+.....+19)]-2[\log1+\log2+\log3+.....+\log7]$, where a, b and c are positive integers. The value of $2a+3b+5c$ is equal to (where $\log a=\log_{10}a$)
\begin{qoptions}
\opt{A}{72}
\opt{B}{42}
\opt{C}{52}
\opt{D}{32}
\end{qoptions}
\qn{3} If the equation $x^4-(7k+11)x^2+25k^2=0$ (k>0) has four real solutions which are in arithmetic progression, then value of $k=\dfrac ab$ where `a' and `b' are relatively prime, then find the value of $a-b$
\begin{qoptions}
\opt{A}{1}
\opt{B}{2}
\opt{C}{4}
\opt{D}{8}
\end{qoptions}
\qn{4} Let $f(x)=2x^2+px+1$ is given. If $f(x)$ is negative integer for only one real value of x, then product of all possible values of p is -
\begin{qoptions}
\opt{A}{$-3$}
\opt{B}{$-16$}
\opt{C}{5}
\opt{D}{$-7$}
\end{qoptions}
\qn{5} Let x,y,z be real numbers such that x, 5y and 7z form a geometric progression while $\dfrac1x,\dfrac1y,\dfrac1z$ form an arithmetic progression. Then the value of $\dfrac7{86}\left(\dfrac xz+\dfrac zx\right)$ is
\begin{qoptions}
\opt{A}{0}
\opt{B}{1}
\opt{C}{2}
\opt{D}{3}
\end{qoptions}
\qn{6} Let p,q and r are all different and non-zero real numbers (taken in that order) are in arithmetic progression. If the roots of quadratic equation $px^2+qx+r=0$ are a and b such that $a^{-1}+b^{-1}$, $a+b$, $a^2+b^2$ (taken in that order) are in geometric progression, then $\dfrac pr$ is equal to-
\begin{qoptions}
\opt{A}{2}
\opt{B}{3}
\opt{C}{$\dfrac23$}
\opt{D}{$\dfrac32$}
\end{qoptions}
\sect{SECTION-I(ii)}
\qn{7} Positive integers a,b and c are chosen so that a < b < c and the system of equations $2x+y=2003$ and $y=|x-a|+|x-b|+|x-c|$ has exactly one solution, then-
\begin{qoptions}
\opt{A}{minimum value of c is 1002}
\opt{B}{maximum value of c is 1002}
\opt{C}{maximum value of b is 1001}
\opt{D}{minimum value of b is 1001}
\end{qoptions}
\qn{8} If a,b,c > 0, such that $a^2+b^2+c^2=27$, then-
\begin{qoptions}
\opt{A}{maximum value of (a + b + c) is equal to 9}
\opt{B}{maximum value of abc is equal to 27}
\opt{C}{minimum value of $\dfrac1a+\dfrac1b+\dfrac1c$ is equal to 1}
\opt{D}{maximum value of $\dfrac1a+\dfrac1b+\dfrac1c$ is equal to 1}
\end{qoptions}
\qn{9} The common difference of the arithmetic progression $(a_n)$ is d and the common ratio of the geometric progression $(b_n)$ is r, where r is a positive rational number less than 1 and $n\in$ N. Given $a_1=d$, $b_1=d^2$ and $E=\dfrac{a_1^2+a_2^2+a_3^2}{b_1+b_2+b_3}$ is a positive integer, then-
\begin{qoptions}
\opt{A}{$\left[\dfrac{E+r}{E-5r}\right]=1$}
\opt{B}{$\left[E-\dfrac12\right]+\left\{r+\dfrac12\right\}=7$}
\opt{C}{$E^2+\left(r+\dfrac32\right)^2=68$}
\opt{D}{$E-2r=9$}
\end{qoptions}
(where [.], \{.\} denotes greatest integer function and fractional part function respectively)
\qn{10} If the cubic equation $x^3-x+1=0$ has roots $\alpha,\beta,\gamma$, then-
\begin{qoptions}
\opt{A}{$\alpha^2+\beta^2+\gamma^2=2$}
\opt{B}{$\alpha^3+\beta^3+\gamma^3=-3$}
\opt{C}{$\alpha^4+\beta^4+\gamma^4=2$}
\opt{D}{$\alpha^5+\beta^5+\gamma^5=-5$}
\end{qoptions}
\qn{11} For all real number x, $2^x+3^x-4^x+6^x-9^x$ is always less than -
\begin{qoptions}
\opt{A}{0}
\opt{B}{1}
\opt{C}{2}
\opt{D}{3}
\end{qoptions}
\qn{12} If the ratio of roots of the equation $x^2+kx+k+2=0$ is 2, the possible value of k is
\begin{qoptions}
\opt{A}{3}
\opt{B}{6}
\opt{C}{$-\dfrac12$}
\opt{D}{$-\dfrac32$}
\end{qoptions}
\sect{SECTION-II}
\qn{1} If $\displaystyle\sum_{r=1}^{44}\log_2(\sin2r^\circ)-\sum_{r=1}^{89}\log_2(\cos r^\circ)=k$, then the value of $\left|\dfrac k{10}\right|$ is
\qn{2} Let a, b, c, d are zeroes of $\alpha(x)=10x^4+px^3+qx^2+rx+s$ and a, c, d are zeroes of $\beta(x)=x^3-9x^2+tx-24$, where p, q, r, s, t $\in$ R and a < b < c < d. If a, c, d (taken in that order) are in arithmetic progression and a, b, c, d (taken in that order) are in harmonic progression and $n=\left(\sqrt{\dfrac{\alpha(6)}{\beta(6)}}-\sqrt{\dfrac{\alpha(5)-\beta(5)}{\beta(5)}}\right)$, then the value of $\dfrac n4$ is
\qn{3} If the number of negative solutions of the equation $x^22^{x+1}+2^{|x-3|+2}=x^2.2^{|x-3|+4}+2^{x-1}$ is k then the value of $\dfrac k2$ is
\qn{4} Let a,b,c be the roots of the equation $x^4+x^3+x^2+x-1=0$. Let $f(x)=x^6-6x^2+6x+7$ and $g(x)=px^2+qx+r$ (p,q,r $\in$ R, p $\ne$ 0). If $f(a)=g(a)$; $f(b)=g(b)$ and $f(c)=g(c)$; then the value of $\dfrac2{g(1)}$ is
\qn{5} A quadratic equation with roots $\alpha$ and $\beta$ where $\alpha\beta=4$, satisfies the equation $\dfrac\alpha{\alpha-1}+\dfrac\beta{\beta-1}=\dfrac{a^2-7}{a^2-4}$ . Number of integral values of `a' for which $\alpha,\beta\in(1,4)$ is equal to
\qn{6} If P(x) is a polynomial degree four with leading coefficient one such that P(1) = 2, P(2) = 6, P(3) = 12, P(4) = 20. Also P(x) = 0 has roots $\alpha,\beta,\gamma$ and $\delta$. Then the value of $\dfrac{(\alpha+3)(\beta+3)(\gamma+3)(\delta+3)}{100}$ is
\qfullbegin
\papersec{2}
\partsec{1}{PHYSICS}
\sect{SECTION-I}
\qn{1} A cylindrical solid pipe of radius $r$ is kept on a truck as shown in figure. If the truck now starts moving with a constant acceleration of $a_0$, the pipe rolls backward without slipping on the floor of the truck and finally falls on the road. If the pipe moves a total length of $l$ on the floor of the truck, then\par
\qfig{p2_phy_s1_q01_1}
\qfullsplit
\begin{qoptions}
\opt{A}{Speed of pipe relative to truck at the time it starts falling is $\dfrac2{\sqrt3}\sqrt{la_0}$\,.}
\opt{B}{Speed of pipe with respect to ground at the time it starts falling is $\dfrac1{\sqrt3}\sqrt{la_0}$\,.}
\opt{C}{Time after which cylinder leaves the truck is $\sqrt{\dfrac{3l}{a_0}}$\,.}
\opt{D}{Angular velocity of the pipe at the time of losing contact with the floor of truck is $\dfrac2{\sqrt3}\dfrac{\sqrt{la_0}}r$\,.}
\end{qoptions}
\qfullend
\qn{2} A non-homogeneous thin rod AB of mass 2M \& length L is kept as shown in figure, with the help of two rigid supports at its end. The mass per unit length of rod is $\lambda=kx^2$ where K is a positive constant and x is distance from end A. Friction is absent everywhere then choose the correct options:\par
\qfig{p2_phy_s1_q02_1}
\begin{qoptions}
\opt{A}{The centre of mass of the rod is at $\dfrac{3L}4$ from the end A.}
\opt{B}{Reaction force due to hinge support at end A \& B are $\dfrac{Mg}2\ \&\ \dfrac{3Mg}2$ respectively.}
\opt{C}{If the hinge at end B is removed suddenly, the angular acceleration of the rod about end A is $\dfrac{5g}{4L}$.}
\opt{D}{After the hinge at B is removed and when the rod becomes vertical, the angular velocity of the rod is $\sqrt{\dfrac{5g}{2L}}$.}
\end{qoptions}
\qn{3} A uniform disc of mass 5 kg and radius 1 m is placed on a rough horizontal surface. A horizontal time varying force is applied on the centre of disc and its variation with time is shown in the graph. The co-efficient of friction between disc and surface is 0.4. (Take g = 10 m/s$^2$)\par
\qfig{p2_phy_s1_q03_2}
\qfig{p2_phy_s1_q03_1}
Then which of the following option/s is/are correct
\begin{qoptions}
\opt{A}{Disc rolls without slipping at t = 3 second.}
\opt{B}{Disc rolls without slipping at t = 4 second.}
\opt{C}{Disc starts slipping at t = 6 second.}
\opt{D}{Friction force is 10 N at t = 3 second.}
\end{qoptions}
\qfullbegin
\qn{4} 5 elastic balls A,B, C, D and E of same mass are moving on a long frictionless horizontal wire. A graph is drawn for their position x against time. Select the possible \textbf{CORRECT} statement(s) on basis of graph.\par
\qfig{p2_phy_s1_q04_1}
\qfullsplit
\begin{qoptions}
\opt{A}{There are a total of 10 collisions.}
\opt{B}{Ball A finally moves with a velocity of +1.5 m/s.}
\opt{C}{Ball C finally moves with a velocity of 0.5 m/s.}
\opt{D}{Ball D finally moves with a velocity of $\dfrac{-10}9$ m / s\,.}
\end{qoptions}
\qfullend
\qn{5} A large wedge BCD having its inclined surface at an angle $\theta$ = 45$^\circ$ to the horizontal, is travelling horizontally leftwards with uniform velocity v$_1$ = 10 m/s. At some instant, a particle is projected vertically up with speed v$_2$ = 20 m/s from point A on ground lying at some distance left to the lower edge B of the wedge as shown in figure. The particle strikes the incline BC normally, while it was falling.\par
\qfig{p2_phy_s1_q05_1}
\begin{qoptions}
\opt{A}{The distance AB at the instant the particle was projected from A, is 15 m.}
\opt{B}{The distance of lower edge B of the wedge from point A at the instant particle strikes the incline, is 15 m.}
\opt{C}{The distance of lower edge B of the wedge from point A at the instant particle strikes the incline, is 30 m.}
\opt{D}{The path of the particle in the reference frame attached to wedge is parabolic.}
\end{qoptions}
\qn{6} Figure shows a frictionless track that consists of a horizontal part of unknown length, which connects to a vertical semicircle of radius r as shown. An object is given an initial velocity v$_0$, (so that it is just able to reach the topmost point) is to move in such a way that after leaving the semicircle at the top it is to fall back to its initial position. Mark the \textbf{CORRECT} statement(s) :\par
\qfig{p2_phy_s1_q06_1}
\begin{qoptions}
\opt{A}{Length AB = 2R}
\opt{B}{v$_0=\sqrt{5\mathrm{gR}}$}
\opt{C}{Tangential acceleration at C is g}
\opt{D}{Normal reaction at C is 3mg}
\end{qoptions}
\sect{SECTION-II}
\qn{1} A particle of mass 0.5 kg is projected with a speed of 20 m/s at an angle of 60$^\circ$ with the horizontal. A bird of mass $\dfrac13$ kg observes the trajectory of the particle and follows exactly the same path with a constant speed of 5 m/s. The force exerted (in N) by air on the bird at the highest point will be (neglect air resistance on particle)
\qn{2} A solid sphere of mass m = 3kg and radius R = 10 cm is projected up an inclined plane of inclination $\theta$ = 30$^\circ$ with an initial velocity v$_0$ = 15 m/s and initial angular velocity $\omega_0=\dfrac{\mathrm v_0}{3\mathrm R}$ as shown. The coefficient of friction between the sphere and inclined plane is $\mu$ = 2tan$\theta$. Then find the velocity (in m/s) of the sphere, when it will start rolling without slipping relative to the inclined plane.\par
\qfig{p2_phy_s2_q02_1}
\qn{3} A thin uniform rod of length 3m is kept vertical at the edge of a tabletop, is disturbed from rest from the vertical position. The lower end of the rod does not slip as the rod turns in the vertical plane shown. Find the angular velocity (in rad/s) of the rod when it leaves contact with the tabletop.\par
\qfig{p2_phy_s2_q03_1}
\qn{4} A box is being pulled by two tractors. One of the tractors has a velocity of v$_1$ = 5 m/s and other has velocity of v$_2$ = 3m/s. The ropes are parallel to velocity vectors. Find the velocity of the box (in m/s). Box is not rotating.\par
\qfig{p2_phy_s2_q04_1}
\qn{5} A smooth pipe placed in a vertical plane is made-up of four quarter circular segments attached as shown in the figure. A rope of length equal to the length of the pipe is passed through the pipe and is hinged at point A. If the initial acceleration of the rope when released from rest by removing hinge at A is Xg, find the value of 2X.\par
\qfig{p2_phy_s2_q05_1}
\qn{6} As shown in the figure below, a rectangular block of mass M and length L is placed on a smooth table. A small cube of mass m with negligible length is placed on the left end of the block with horizontal initial velocity v sliding to the right. The coefficient of kinetic friction between the block and the cube is $\mu$, and the cube does not fall from the right end of the block. What is the displacement (in m) of the rectangular block when the cube stops sliding on the rectangular block exactly at its right most end ? (Take m = 2kg , M = 6kg and L = 10m)\par
\qfig{p2_phy_s2_q06_1}
\sect{SECTION-III}
\qn{1} A bead can slide on smooth horioztnal fixed frame of shape of quater circle of radius R $=\sqrt2$m with C as center of curvature. A force F = K(3 $-$ y) is applied on the bead always towards point B, where y (in metre) is separation between position of bead and point B and K = 1 N/m. Find work done by this force from A to B in J.\par
\qfig{p2_phy_s3_q01_1}
\qn{2} Three identical spherical balls kept in equilibrium inside a hemispherical bowl at its bottom. The angle between line of N$_1$ and line of N$_2$, if radius of balls and bowl adjusted such that normal force between B and C approaching to zero, is $\tan^{-1}\left(\dfrac{\sqrt{\mathrm x}}{\mathrm y}\right)$. (x and y are minimum possible integer; N$_1\to$ normal between A and B; N$_2\to$ Normal between B and bowl and all surfaces are smooth). Find the (x + y).\par
\qfig{p2_phy_s3_q02_1}
\qn{3} Through a smooth fixed cylinder of radius R, the axis of which is horizontal, a weightless rope is thrown, to the ends of which a load and a thin uniform rod is attached as shown in the figure. In the position of equilibrium, rod makes an angle $\alpha$ = 30$^\circ$ with the horizontal. The distance from the end of rod to which the rope is attached, to the point of contact of the rod and the cylinder is $\dfrac R{\sqrt2}$. If mass M of rod is 6 kg, then find the value of mass m in kg.\par
\qfig{p2_phy_s3_q03_1}
\qn{4} A uniform solid sphere rolling without slipping along a horizontal plane with angular speed $\omega_0$ = 14 rad/sec, suddenly collides with a plane inclined at an angle $\theta=60^\circ$ as shown in the figure. Find the angular velocity (in rad/sec) of sphere just after impact : (assume that there is no slipping anytime and collision is perfectly inelastic)\par
\qfig{p2_phy_s3_q04_1}
\qn{5} At an instant t = 0, Block A of mass 1 kg is moving with speed 8 m/s towards right on rough surface of block B of mass 3 kg. Block B, which is placed on smooth horizontal surface is moving with speed 4 m/s towards right at same instant (t = 0). The net work done (in J) by the frictional force in long time is (block A will not fall from block B)\par
\qfig{p2_phy_s3_q05_1}
\qn{6} A water channel of width b = 17 m lies between two frictionless platforms A and B. A plank (P) of length $l$ = 2.0 m rest near the platform A. A block of mass m = 25 kg sliding on the platform A lands on the plank and stays on it. Force $\vec{\mathrm F}$ of water resistance on the plank is proportional to velocity of the plank ($\vec{\mathrm v}_{\mathrm{pr}}$) relative to the water and given by the law $\vec{\mathrm F}=-\mathrm k\vec{\mathrm v}_{\mathrm{pr}}$, where k = 15 N-s/m. Calculate the minimum speed u$_{\mathrm{min}}$ of the block on platform A so that plank reaches the platform B. (Answer in m/s)\par
\qfig{p2_phy_s3_q06_1}
\partsec{2}{CHEMISTRY}
\sect{SECTION-I}
\qn{1} Following statement(s) is/are true for D-glucose :
\begin{qoptions}
\opt{A}{It shows Schiff's reagent test}
\opt{B}{It is C-2 epimer of D-mannose}
\opt{C}{Oxidise with Br$_2$ / H$_2$O to form saccharic acid}
\opt{D}{Shows mutarotation.}
\end{qoptions}
\qn{2} Optically active amino acid is :
\begin{qoptions}
\opt{A}{Glycine}
\opt{B}{Alanine}
\opt{C}{Valine}
\opt{D}{Phenyl alanine}
\end{qoptions}
\qn{3} Choose \textbf{CORRECT} statement is/are
\begin{qoptions}
\opt{A}{The oxygen atoms in XeOF$_3^{\oplus}$ and XeO$_2$F$_2$ prefer equitorial positions}
\opt{B}{CO+H$_2$O $\xrightarrow{\text{Catalyst}}$ CO$_2$ + H$_2$; is known as water gas shift reaction}
\opt{C}{KO$_2$ is used as an auxillary supply of O$_2$ in mines, submarines and space vehicles.}
\opt{D}{The superoxides MO$_2$ contains the paramagnetic O$_2^{\ominus}$ which is stable in presence of large cations such as K, Rb, Cs.}
\end{qoptions}
\qn{4} Which of the following statement is \textbf{CORRECT}?
\begin{qoptions}
\opt{A}{Among C$_2$H$_6$, N$_2$H$_4$, H$_2$O$_2$ and F$_2$ strongest bond between same atom is in C$_2$H$_6$.}
\opt{B}{Among Li, Na, K, Rb, and Cs, only Li can react with nitrogen when burn in air.}
\opt{C}{Among the oxide (M$_2$O), peroxide (M$_2$O$_2$) and superoxide (MO$_2$) only superoxide evolve oxygen when reacts with ice cold water.}
\opt{D}{Iodine has a very low solubility in water but dissolves readily in KI (aq.)}
\end{qoptions}
\qn{5} For the reaction A(g) $\rightleftharpoons$ B(g) following graph is obtained -\par
Hence we may conclude for the reaction -\par
\qfig{p2_chem_s1_q05_1}
\begin{qoptions}
\opt{A}{$\Delta G^o<0$ and $K^0_p>1$}
\opt{B}{$\Delta G^o<0$ and $K^0_p<1$}
\opt{C}{$\Delta S^o>0$ and $\Delta G^0>0$}
\opt{D}{$\Delta G^o>0$ and $\Delta H^0>0$}
\end{qoptions}
\qn{6} From the following data, select the correct option(s)\par
Given : $(\Delta_fH^0)_{\mathrm{AlF_3(s)}}=-1500$ kJ/mol , $(\Delta_fH^0)_{\mathrm{Al^{+3}(aq.)}}=-530$kJ/mol\par
$(\Delta_fH^0)_{\mathrm{F^-(aq.)}}=-330$ kJ/mol , $(\Delta_fH^0)_{\mathrm{Al^{+3}(g)}}=+5480$ kJ/mol
\begin{qoptions}
\opt{A}{$(\Delta H_{\mathrm{hydration}})_{\mathrm{Al^{+3}(g)}}=-6010$ kJ/mol}
\opt{B}{$(\Delta H_{\mathrm{solution}})_{\mathrm{AlF_3(s)}}=-20$ kJ/mol}
\opt{C}{$(\Delta H_{\mathrm{solution}})_{\mathrm{AlF_3(s)}}=+20$ kJ/mol}
\opt{D}{$(\Delta H_{\mathrm{hydration}})_{\mathrm{Al^{+3}(g)}}=-4950$ kJ/mol}
\end{qoptions}
\sect{SECTION-II}
\qn{1} If $\epsilon_{\mathrm{H-Br}}=368$ kJ/mol, $\epsilon_{\mathrm{H-H}}=435$ kJ/mol , $\epsilon_{\mathrm{Br-Br}}=192$ kJ/mol , then find magnitude of $\Delta_rH^0$(kJ / mol) for the reaction :\par
$\dfrac12$H$_2$(g) $+$ $\dfrac12$Br$_2$(g) $\to$ HBr(g)
\qn{2} If $(\Delta_fH^\circ)_{\mathrm{H^+(g)}}=1536$ kJ/mol and enthalpy of atomisation of H$_2$(g) is 436 kJ/mol. then ionisation enthalpy of H atom in kJ/mol is :
\qn{3} Total number of oxyacid(s) which have basicity more than two and contain atleast one p$\pi$--p$\pi$ bond between same atoms\par
H$_4$P$_2$O$_5$, H$_2$S$_2$O$_8$, H$_2$S$_2$O$_3$, H$_3$BO$_3$, H$_2$N$_2$O$_2$, H$_4$SiO$_4$, H$_4$P$_2$O$_7$, H$_6$Si$_2$O$_7$
\qn{4} Consider the following reaction for the production of UF$_6$ in the enrichment of uranium (U$^{235}$)\par
U$_{\mathrm{(s)}}$ + interhalogen$_{\mathrm{(A)}}$ $\longrightarrow$ UF$_6$ + interhalogen$_{\mathrm{(B)}}$\par
A and B are interhalogen compound made by chlorine and fluorine, then calculate the value of $\left(\dfrac xy\right)$ where x is total number of lone pairs present in one molecule of (A) and y is total number of fluorine atom/s present in one molecule of (B).
\qfullbegin
\qn{5} Number of reaction(s) which are correctly represented?\par
\qfig{p2_chem_s2_q05_1}
\qfullend
\qn{6} Count number of polymers which are condensation polymers as well as co-polymers :\par
(1) Nylon-6 \quad (2) Nylon-6, 6 \quad (3) Nylon-6, 10 \quad (4) Dacron\par
(5) Bakelite \quad (6) Melamine formaldehyde resin \quad (7) Polyisoprene\par
(8) Styrene-butadiene rubber (Buna-S) \quad (9) PAN \quad (10) Teflon
\sect{SECTION-III}
\qn{1} 2 mol He at 300 K and 2 atm undergoes sudden adiabatic expansion against 1 atm pressure If decrease in enthalpy of the system is x cal then x/100 is.
\qn{2} Pressure on 1 kg liquid H$_2$O at 4$^\circ$C and 1 bar is increased to 11 bar isothermally then change in Gibb's energy (in kJ) for the process is.
\qfullbegin
\qn{3} The number of reducing sugars among the following is \underline{\hspace{0.8cm}}\par
\qfig{p2_chem_s3_q03_1}
\qn{4} Among the following, the number of aromatic compound(s) is \underline{\hspace{0.8cm}}\par
\qfig{p2_chem_s3_q04_1}
\qfullend
\qn{5} Total number of reaction(s) in which oxygen is evolved.\par
(A) H$_2$O$_2\xrightarrow{\Delta}$ \quad (B) H$_2$O$_2$ + Fe$^{3\oplus}\xrightarrow{\text{basic medium}}$ \quad (C) H$_2$O$_2$ + Fe$^{2\oplus}\xrightarrow{\text{acidic medium}}$\par
(D) H$_2$O$_2$ + SO$_3^{2-}\longrightarrow$ \quad (E) H$_2$O$_2$ + MnO$_4^{\ominus}\xrightarrow{\mathrm{H^{\oplus}}}$ \quad (F) H$_2$O$_2$ + O$_3\longrightarrow$\par
(G) H$_2$O$_2$ + HClO $\longrightarrow$
\qn{6} Total number of specie(s) which give hydrogen gas with dil. HCl\par
Cu, Ag, Cr$^{2\oplus}$, Zn, Al, Cr, Na
\partsec{3}{MATHEMATICS}
\sect{SECTION-I}
\qn{1} Which of the following statement(s) is/are true ?
\begin{qoptions}
\opt{A}{$\log_{10}3$ lies between $\dfrac13$ and $\dfrac12$}
\opt{B}{$\log_{\mathrm{cosec}\frac\pi3}\cos\dfrac\pi6=1$}
\opt{C}{$e^{\ell n(\ell n7)}$ is greater than 1}
\opt{D}{$\log_{10}2+\dfrac12\log_{10}5+\log_{10}\left(2+\sqrt5\right)=\log_{10}\left(4\sqrt5+10\right)$}
\end{qoptions}
\qn{2} If $\alpha$ and $\beta$ are the roots of the equation $x^2-6x-4=0$, then $\dfrac{\alpha^{17}+\beta^{17}-4\left(\alpha^{15}+\beta^{15}\right)}{\left(\alpha^8+\beta^8\right)^2-2^{17}}$ is a-
\begin{qoptions}
\opt{A}{real number}
\opt{B}{composite number}
\opt{C}{prime number}
\opt{D}{rational number}
\end{qoptions}
\qn{3} The equation $\sqrt{x+3-4\sqrt{x-1}}+\sqrt{(x+15)-8\sqrt{x-1}}=2$ has
\begin{qoptions}
\opt{A}{no real root}
\opt{B}{exactly one real root}
\opt{C}{exactly two real roots}
\opt{D}{infinite number of real roots}
\end{qoptions}
\qn{4} If $x_1,x_2$ are real solutions to the equation $x^2+a^2x+b=0$ while $x_3,x_4$ are real solutions to $x^2+5ax+7=0$ such that $x_1-x_3=x_2-x_4=2$, then-
\begin{qoptions}
\opt{A}{$x_1,x_2,x_3,x_4\in Q$}
\opt{B}{$x_1,x_2,x_3,x_4\notin Q$}
\opt{C}{$a+b=2$}
\opt{D}{$a+b=-25$}
\end{qoptions}
\qn{5} Let $x=\underbrace{111...11}_{\text{20 digits}}$, $y=\underbrace{444...44}_{\text{20 digits}}$, $z=\underbrace{777...77}_{\text{10 digits}}$, $w=\underbrace{999...99}_{\text{10 digits}}$ then choose the correct option(s).
\begin{qoptions}
\opt{A}{$\dfrac{49x-9z^2}{14w}=\dfrac79$}
\opt{B}{$\dfrac{49x-9z^2}{14w}=\dfrac29$}
\opt{C}{$9y=4w^2+8w$}
\opt{D}{$9y=6w^2+8w$}
\end{qoptions}
\qn{6} Let $f(x)=x^{2018}+a_0x^{2017}+a_1x^{2016}+a_2x^{2015}+.....+a_{2017}$ where $a_0,a_1,a_2,......,a_{2017}$ are real numbers such that $f(0)=f(1)$ and each root of $f(x)=0$ is real and always lies between 0 and 1 then product of roots is
\begin{qoptions}
\opt{A}{less than or equal to $\dfrac1{2^{2018}}$}
\opt{B}{less than or equal to $\dfrac1{2^{2017}}$}
\opt{C}{greater than or equal to $\dfrac1{2^{2018}}$}
\opt{D}{greater than $\dfrac1{2^{2019}}$}
\end{qoptions}
\sect{SECTION-II}
\qn{1} If $a^2+8b^2+18c^2-4ab-12bc-36c+36=0$, where a, b, c $\in$ R, then the value of $\dfrac a{4bc}$ is
\qn{2} If sum of all the integers in the range of $f(x)=\sqrt{21-\left(\sum_{k=1}^{7}|x-k|\right)}$ is `$\mu$' then digit $\dfrac\mu8$ equals
\qn{3} If $a_1,a_2,.......,a_{10}$ are positive numbers in arithmetic progression such that\par
$\dfrac1{a_1a_2}+\dfrac1{a_2a_3}+.....+\dfrac1{a_9a_{10}}=\dfrac9{64}$ and $\dfrac1{a_1a_{10}}+\dfrac1{a_2a_9}+......\dfrac1{a_{10}a_1}=\dfrac1{10}\left(\dfrac1{a_1}+....+\dfrac1{a_{10}}\right)$,\par
then the value of $\dfrac25\left(\dfrac{a_1}{a_{10}}+\dfrac{a_{10}}{a_1}\right)$ is
\qn{4} If the value of $\displaystyle\sum_{n=1}^{2016}\dfrac{\sqrt{n^4+2n^3+3n^2+2n+1}}{n(n+1)}$ is expressed in the form of $\lambda+\dfrac\lambda{\lambda+1}$ then sum of digits of $\lambda$ equals
\qn{5} If largest real root of the equation $x^4-4x^3+5x^2-4x+1=0$ can be expressed as $\dfrac{a+\sqrt b}c$ (where does not contain any perfect square and a, b, c are natural numbers in lowest form), then the value of $\left(\dfrac{a+b-c}4\right)$ is
\qn{6} If the ordered pair (a,b) satisfying the equation $a^2=1+6\log_4b$ and $b^2=2^ab+2^{2a+1}$ is $(a_1,b_1)$ and $(a_2,b_2)$ then value of $\log_4|a_1a_2b_1b_2|$ is
\sect{SECTION-III}
\qn{1} If $\dfrac{\log_{10}x+\log_{10}y}{(a+b)}=\dfrac{\log_{10}y+\log_{10}z}{(b+c)}=\dfrac{\log_{10}z+\log_{10}x}{(c+a)}$, then $x^{(c-b)}y^{(a-c)}z^{(b-a)}$ is (wherever defined)
\qn{2} Number of real ordered pairs (x, y) satisfying $4^{\frac xy+\frac yx}=32$ and $\log_3(x-y)+\log_3(x+y)=1$, is
\qn{3} If K is the sum of all possible integral value(s) of $n\in[1,300]$ for which the quadratic equation $x^2-3x-n=0$ has integral roots then $\dfrac K{400}$ is
\qn{4} Number of integer(s) satisfying the equation $\left(\left|x^2-5x+5\right|\right)^{x^2-9x+20}=1$, is equal to
\qn{5} Number of integral values of `a' for which the equation $(x^2+x+2)^2-(a-3)(x^2+x+2)(x^2+x+1)+(a-4)(x^2+x+1)^2=0$ has atleast one real root is
\qn{6} If a, b, c are three distinct positive numbers are in H.P. such that $\left(\dfrac{a+b}{2a-b}\right)+\left(\dfrac{c+b}{2c-b}\right)>\sqrt{\lambda\sqrt{\lambda\sqrt{\lambda\sqrt{\lambda........\infty}}}}$, then maximum integral value of $\lambda$ is equal to
\end{multicols}
\end{document}
